\documentclass[tikz,border=6pt]{standalone}
\usepackage[no-math]{fontspec}
\usepackage[brazil]{babel}
\usepackage{amsmath,amssymb}
\IfFileExists{C:/Windows/Fonts/Lato_Regular.ttf}{%
  \setmainfont{Lato}[
    Path = C:/Windows/Fonts/,
    UprightFont = Lato_Regular.ttf,
    BoldFont = Lato_Bold.ttf,
    ItalicFont = Lato_Italic.ttf,
    BoldItalicFont = Lato_BoldItalic.ttf
  ]
  \setsansfont{Lato}[
    Path = C:/Windows/Fonts/,
    UprightFont = Lato_Regular.ttf,
    BoldFont = Lato_Bold.ttf,
    ItalicFont = Lato_Italic.ttf,
    BoldItalicFont = Lato_BoldItalic.ttf
  ]
}{%
  \RequirePackage[default]{lato}
}
\renewcommand{\familydefault}{\sfdefault}

\usepackage{xcolor}
\usepackage{tikz}
\usepackage{pgfplots}
\pgfplotsset{compat=1.18}
\usetikzlibrary{arrows.meta,positioning,shapes.geometric,fit,backgrounds,calc,decorations.pathreplacing,patterns}

% Paleta institucional UEFS / CHF662
\definecolor{uefsNavy}{HTML}{1B2A4A}
\definecolor{uefsTeal}{HTML}{0E7C7B}
\definecolor{uefsForest}{HTML}{2D6A4F}
\definecolor{uefsGold}{HTML}{C4983B}
\definecolor{uefsGoldDark}{HTML}{8F6B1E}
\definecolor{uefsAlert}{HTML}{B02A37}
\definecolor{uefsOrange}{HTML}{D95F02}
\definecolor{uefsInk}{HTML}{22252A}
\definecolor{uefsMuted}{HTML}{5C6470}
\definecolor{uefsWater}{HTML}{2E86AB}
\definecolor{uefsSand}{HTML}{D9B382}
\definecolor{bgNavy}{HTML}{F2F5F9}
\definecolor{bgTeal}{HTML}{E9F5F4}
\definecolor{bgForest}{HTML}{EEF5F0}
\definecolor{bgGold}{HTML}{FCF7EB}
\definecolor{bgAlert}{HTML}{FDF2F2}
\definecolor{bgGray}{HTML}{F8FAFC}
\definecolor{lineGray}{HTML}{CFD8DC}

\tikzset{
  >=Stealth,
  aaCard/.style={draw=#1, fill=white, rounded corners=4pt, line width=1pt, inner sep=6pt, align=center, font=\small},
  aaBox/.style={draw=#1, rounded corners=4pt, line width=1.1pt, inner sep=6pt, align=center, font=\footnotesize\sffamily, text=uefsInk},
  aaTitle/.style={font=\normalsize\bfseries\sffamily, text=uefsNavy, align=center},
  aaSub/.style={font=\footnotesize\sffamily, text=uefsMuted, align=center},
  aaArrow/.style={->, >=Stealth, line width=1.4pt},
  aaDArrow/.style={<->, >=Stealth, line width=1.4pt}
}


\begin{document}
\begin{tikzpicture}[font=\sffamily]

% Macro para árvore na encosta
\newcommand{\arvoreEnc}[2]{%
  \draw[brown!65!black, line width=2.2pt] (#1,#2) -- (#1,#2+0.38);
  \fill[uefsForest!85] (#1,#2+0.68) ellipse (0.33cm and 0.39cm);
  \draw[uefsForest!95!black, line width=0.6pt] (#1,#2+0.68) ellipse (0.33cm and 0.39cm);
}

% =========================================================
% PAINEL (a) Vertente 1: Parque Estadual Preservado
% =========================================================
\begin{scope}[shift={(0,0)}]
  \draw[draw=uefsForest, fill=bgGray, rounded corners=6pt, line width=1.2pt] (0,0) rectangle (9.6, 7.6);
  \node[font=\normalsize\bfseries, text=uefsForest] at (4.8, 7.05) {(a) Vertente 1: Parque Estadual Preservado};

  % Corpo da encosta (topo -> ombro -> vertente 38% -> sopé -> planície)
  \fill[uefsNavy!13]
    (0.4, 5.45) -- (2.4, 5.2) -- (6.6, 2.35) -- (9.2, 2.1) -- (9.2, 1.55) -- (0.4, 1.55) -- cycle;

  % Linha de superfície topográfica
  \draw[line width=1.6pt, uefsNavy] (0.4, 5.45) -- (2.4, 5.2) -- (6.6, 2.35) -- (9.2, 2.1);

  % Vegetação nativa preservada
  \arvoreEnc{1.1}{5.36}
  \arvoreEnc{2.25}{5.22}
  \arvoreEnc{3.4}{4.52}
  \arvoreEnc{4.55}{3.74}
  \arvoreEnc{5.65}{3.0}

  % Rótulos no interior amplo do maciço (longe das bordas e da linha de encosta)
  \node[align=center, font=\footnotesize\bfseries, text=uefsInk] at (2.8, 2.55)
    {Neossolo Litólico raso\\Declividade = 38\%};

  \node[align=center, font=\footnotesize\bfseries, text=uefsForest] at (7.8, 3.15)
    {Sem ocupação\\humana na base};

  % Quadro inferior de síntese
  \node[aaBox=uefsForest, fill=white, text width=8.4cm, minimum height=1.1cm, inner sep=5pt]
    at (4.8, 0.8) {
      \footnotesize\textbf{Alta Fragilidade Física (38\%)}\\
      Perigo Presente $\vert$ Vulnerabilidade Nula ($V = 0$)\\
      \textbf{\color{uefsForest}Risco Humano: $R = P \times V = 0$}
    };
\end{scope}

% =========================================================
% PAINEL (b) Vertente 2: Encosta Desmatada e Ocupada
% =========================================================
\begin{scope}[shift={(10.2, 0)}]
  \draw[draw=uefsOrange, fill=bgGray, rounded corners=6pt, line width=1.2pt] (0,0) rectangle (9.6, 7.6);
  \node[font=\normalsize\bfseries, text=uefsOrange!90!black] at (4.8, 7.05) {(b) Vertente 2: Encosta Desmatada e Ocupada};

  % Corpo da encosta
  \fill[uefsNavy!13]
    (0.4, 5.45) -- (2.4, 5.2) -- (6.6, 2.35) -- (9.2, 2.1) -- (9.2, 1.55) -- (0.4, 1.55) -- cycle;

  % Linha de superfície topográfica
  \draw[line width=1.6pt, uefsNavy] (0.4, 5.45) -- (2.4, 5.2) -- (6.6, 2.35) -- (9.2, 2.1);

  % Vetor de escorregamento / erosão em encosta desmatada (paralelo à vertente, deslocado acima)
  \draw[->, >=Stealth, dashed, line width=2pt, uefsOrange] (3.1, 4.92) -- (6.1, 2.88);
  % Símbolos de blocos/instabilidade
  \fill[uefsOrange] (3.55, 4.62) -- (3.85, 4.42) -- (3.85, 4.82) -- cycle;
  \fill[uefsOrange] (4.75, 3.80) -- (5.05, 3.60) -- (5.05, 4.00) -- cycle;
  \node[font=\scriptsize\bfseries, text=uefsAlert, align=center] at (5.3, 4.65)
    {Instabilidade e\\fluxo de detritos};

  % Rótulo no interior amplo do maciço
  \node[align=center, font=\footnotesize\bfseries, text=uefsInk] at (2.8, 2.55)
    {Neossolo Litólico raso\\Declividade = 38\%};

  % Duas moradias na base (alinhadas ao sopé y~2.2)
  \foreach \hx/\hy in {6.85/2.29, 7.95/2.19} {
    \fill[uefsAlert!30] (\hx, \hy) rectangle (\hx+0.78, \hy+0.55);
    \draw[uefsOrange!80!black, line width=1pt] (\hx, \hy) rectangle (\hx+0.78, \hy+0.55);
    \fill[uefsOrange!90!black] (\hx-0.08, \hy+0.55) -- (\hx+0.39, \hy+0.92) -- (\hx+0.86, \hy+0.55) -- cycle;
  }
  \node[font=\footnotesize\bfseries, text=uefsOrange!90!black] at (7.8, 3.55) {Moradias na base};

  % Quadro inferior de síntese
  \node[aaBox=uefsOrange, fill=white, text width=8.4cm, minimum height=1.1cm, inner sep=5pt]
    at (4.8, 0.8) {
      \footnotesize\textbf{Alta Fragilidade Física (38\%)}\\
      Perigo Crítico ($P = 4$) $\vert$ Vulnerabilidade ($V = 4$)\\
      \textbf{\color{uefsAlert}Risco Crítico: $R = P \times V = 16$}
    };
\end{scope}

\end{tikzpicture}
\end{document}
